\documentclass{article}
\usepackage{amsmath, amssymb, amsthm, mathrsfs}
\usepackage[utf8]{inputenc}
\usepackage[T1]{fontenc}
\usepackage{xcolor}

% Theorem environments
\newtheorem{theorem}{Theorem}[section]
\newtheorem{lemma}[theorem]{Lemma}
\newtheorem{definition}[theorem]{Definition}
\newtheorem{proposition}[theorem]{Proposition}

% Macros for Lean code references
\newcommand{\lean}[1]{\texttt{#1}}
\newcommand{\mathlib}[1]{\texttt{#1}}

% Comparison Annotation Macro
\newcommand{\annotation}[1]{
    \vspace{0.3cm}
    \noindent\fbox{\parbox{\textwidth}{\small \textbf{Comparison with PDF [arXiv:2503.07846]:} #1}}
    \vspace{0.3cm}
}

% Mathematical notation
\newcommand{\adjoin}[2]{#1[#2]}
\newcommand{\minpoly}{\mathrm{minpoly}}
\newcommand{\mk}{\mathrm{mk}}
\newcommand{\aeval}{\mathrm{aeval}}
\newcommand{\residue}{\mathrm{residue}}

\title{Annotated Deformalization: Monogenic Extensions}
\author{Grant Yang (Deformalized \& Annotated by AI)}
\date{\today}

\begin{document}
\maketitle
\section{Introduction}

This document translates the Lean formalization of Monogenic Extensions and compares it directly to Lemmas 3.1 and 3.2 of the source paper \textit{Number Fields Generated by Points in Linear Systems on Curves}[cite: 1].


\section{Basic Definitions and Lemma 3.2 Equivalents}

\begin{lemma}[\lean{monogenic\_of\_univQuot}]
If there exists a polynomial $f \in R[X]$ and an isomorphism $e: R[X]/\langle f \rangle \simeq_a[R] S$, then there exists $\beta \in S$ such that $\adjoin{R}{\{\beta\}} = \top$.
\end{lemma}

\begin{lemma}[\lean{deriv\_isUnit\_of\_monogenic}]
Assume the extension is Étale, $R, S$ are domains, and $R$ is integrally closed. Let $\beta$ generate $S$. Then $\aeval(\beta, \minpoly_R(\beta)') \in S^{\times}$.
\end{lemma}

\annotation{
    \textbf{Splitting of Lemma 3.2:} The source paper combines these concepts into a single \textbf{Lemma 3.2}[cite: 178].
    \begin{itemize}
        \item The paper's Lemma 3.2 states that if $R \to S$ is étale, then $S$ is monogenic \textit{and} $f'(\beta)$ is a unit[cite: 178].
        \item The Lean formalization decouples these. \lean{monogenic\_of\_finiteInjectiveEtale} (proved later) establishes the generator existence, while \lean{deriv\_isUnit\_of\_monogenic} establishes the unit derivative property separately.
        \item \textbf{Proof Strategy:} The paper proves $f'(\beta)$ is a unit by reducing to the residue field, showing the reduced polynomial is separable, and lifting the non-zero residue back to a unit in the local ring [cite: 181-184]. The Lean proof \lean{deriv\_isUnit\_of\_monogenic} follows this exactly, utilizing the fact that the residue extension is separable.
    \end{itemize}
}

\section{Nakayama Helpers (The "Case 1" Logic)}

\begin{lemma}[\lean{adjoin\_eq\_top\_of\_quotient\_gen}]
Assume local rings. Let $q \subset S$ be prime. Let $\beta \in S$.
Assume:
1. The image of $\beta$ generates $S/q$.
2. There exists $\pi \in \adjoin{R}{\{\beta\}}$ such that $\mathfrak{m}_S = \langle \pi \rangle \sqcup \mathfrak{m}_R S$.
Then $\adjoin{R}{\{\beta\}} = \top$.
\end{lemma}

\annotation{
    \textbf{Formalizing the Exact Sequence:} This lemma is the formal verification of the "exact sequence" argument used in the proof of \textbf{Lemma 3.1} in the paper [cite: 190-198].
    \begin{itemize}
        \item The paper argues that if $f_1(\beta) \pmod{\mathfrak{m}_R}$ generates $\mathfrak{m}_S/\mathfrak{m}_R S$, then the map from the polynomial ring is surjective via Nakayama[cite: 198].
        \item The Lean code isolates this logic into a standalone lemma. It essentially proves that if you can generate the "vertical" direction (the quotient $S/q$) and the "horizontal" direction (the maximal ideal $\mathfrak{m}_S$), you generate the whole ring.
    \end{itemize}
}

\section{Main Theorem: Lemma 3.1 Equivalent}

\begin{theorem}[\lean{monogenic\_of\_etale\_height\_one\_quotient}]
Assume $R, S$ are domains, $R$ integrally closed, $S$ a Unique Factorization Monoid (UFD). Assume $S$ is finite and faithful over $R$.
Let $q$ be a prime ideal of height 1 in $S$.
If the quotient map $R/(q \cap R) \to S/q$ is Étale, then $S$ is monogenic.
\end{theorem}

\annotation{
    \textbf{Generalization of Assumptions:}
    \begin{itemize}
        \item The paper's \textbf{Lemma 3.1} assumes $R$ and $S$ are \textit{Regular Local Rings}[cite: 174].
        \item The Lean formalization weakens this assumption to $S$ being a \textit{Unique Factorization Monoid (UFD)} (plus Integrally Closed domains).
        \item This is a valid generalization because every Regular Local Ring is a UFD (Auslander-Buchsbaum theorem). The specific property required by the proof is that the height 1 prime $\mathfrak{q}_0$ is principal[cite: 203], which holds in a UFD.
    \end{itemize}
}

\begin{proof}
Let $\varphi = \text{algebraMap } R \ S$. Let $p = q^c$. Let $R_0 = R/p, S_0 = S/q$.
Let $\varphi_0$ be the induced map.

\textbf{Step 1 (Quotient Properties):}
$\varphi_0$ is finite and injective. Since it is Étale by hypothesis, we find a generator $B_0 \in S_0$ for $S_0$ over $R_0$. Let $f_0 = \minpoly_{R_0}(B_0)$.

\textbf{Step 2 (Lifting):}
Lift $B_0$ to $B \in S$. Lift $f_0$ to a monic $f_1 \in R[X]$.

\textbf{Step 3 (Maximal Ideal Decomposition):}
We utilize the property that $\varphi_0$ is Étale, hence formally unramified. Thus $\mathfrak{m}_S = q \sqcup \mathfrak{m}_R S$.

\annotation{
    \textbf{Matches Paper Step:} This corresponds to the paper's deduction: "Since $R_0 \to S_0$ is étale, $\mathfrak{m}_{S_0} = \mathfrak{m}_{R_0} S_0$... i.e., $\mathfrak{m}_S = \mathfrak{q}_0 + \mathfrak{m}_R S$"[cite: 187].
}

\textbf{Step 4 (Principal Ideal):}
Since $S$ is a UFD and $q$ has height 1, $q$ is principal. Let $q = \langle q_0 \rangle$.

\annotation{
    \textbf{Matches Paper Step:} The paper states "Since $S$ is a regular local ring hence a unique factorization domain, so there exists a $q_0 \in S$ that generates the height 1 ideal $\mathfrak{q}_0$" [cite: 200-203].
}

\textbf{Step 5 (Case Analysis):}
Let $y = \aeval(B, f_1)$.
\textit{Case A:} $y \in \mathfrak{m}_S$ and $\langle y \rangle \sqcup \mathfrak{m}_R S = \mathfrak{m}_S$.
We satisfy the conditions of \lean{adjoin\_eq\_top\_of\_quotient\_gen}. Thus $\adjoin{R}{\{B\}} = \top$.

\textit{Case B:} The condition fails.
This implies $y = q_0 \cdot a$ with $a \in \mathfrak{m}_S$.
We establish that $\aeval(B, f_1')$ is a unit.

\annotation{
    \textbf{Matches Paper Step:} This mirrors the paper's check: "If $f_1(\beta) \pmod{\mathfrak{m}_R}$ generates $\mathfrak{m}_S/\mathfrak{m}_R$... we have the desired equality. Assume that $f_1(\beta) \dots$ does not generate" [cite: 198-199].
}

\textbf{Step 6 (Taylor Correction):}
We choose a new generator $B' = B + q_0$.
We prove a Taylor expansion lemma inline:
\[ f(x+h) = f(x) + f'(x)h + h^2 c \]
Applying this to $B'$:
\[ \aeval(B', f_1) = q_0 (a + \aeval(B, f_1') + q_0 c) \]
The term in the parenthesis is a unit, because $\aeval(B, f_1')$ is a unit, while $a \in \mathfrak{m}_S$ and $q_0 \in \mathfrak{m}_S$.
Therefore, $\langle \aeval(B', f_1) \rangle = \langle q_0 \rangle = q$.
We conclude $\adjoin{R}{\{B'\}} = \top$.

\annotation{
    \textbf{High-Level Alignment:}
    The "Taylor Correction" is the core of the proof in both versions.
    \begin{itemize}
        \item \textbf{Paper:} Explicitly calculates $f_1(\beta + q_0) = f_1(\beta) + f_1'(\beta)q_0 + q_0^2 b$[cite: 205].
        \item \textbf{Lean:} Performs the same calculation (formalized as a Taylor remainder lemma) to show that the new generator's evaluation generates the ideal $\mathfrak{q}_0$, satisfying the "Case A" condition for the new element.
        \item \textbf{Simplification in Paper:} The paper assumes $S$ is regular to get $q_0$. The formalization's use of UFD is a strictly stronger result (as it applies to more rings), but functionally the proofs are identical in flow.
    \end{itemize}
}

\end{proof}

\end{document}